\section{Proofs for structured argumentation}
\label{sec:annex-p18-structured-argumentation}

The construction follows the identified tutorial presentation
\cite[Sections 3 and 5]{ModgilPrakken2014ASPICTutorial}.
The arguments below prove the precise statements in the body, including
the typed preference convention and the complete finite examples.

\begin{annexitem}{annex:p18.aspictree}{thm:p18-aspic-tree-properties}{Tree invariants}
\begin{proof}
Give a premise leaf or a nullary rule node height zero. Give a rule node
with children height one plus the maximum of their heights. The class of
finite labelled ordered trees is a set: each tree can be encoded by a finite
prefix-closed set of finite sequences of natural numbers, with labels in the
set consisting of \(K\) and the rule objects. Taking the subset whose labels
and children satisfy the construction gives \(\mathcal A\).
Equivalently, form arguments of height at most \(j\) by induction and take
their union over \(j<\omega\). Each rule has finite arity, so these two
descriptions agree. No infinite branch is admitted.

All the component sets at a leaf or nullary node are finite by their
definitions. At a non-nullary node they are finite unions of finite child
component sets, with at most one additional rule or tree. Induction on height
proves their finiteness.

For subargument transitivity, suppose \(B\in\operatorname{Sub}(A)\).
If \(B=A\), there is nothing to prove. Otherwise \(B\) occurs below some
child \(A_i\). Induction gives
\(\operatorname{Sub}(B)\subseteq\operatorname{Sub}(A_i)
\subseteq\operatorname{Sub}(A)\).
The assertion also covers repeated occurrences of the same subtree:
component sets retain the subtree once, while the argument itself retains
the full ordered tree.

Finally, a premise leaf and a nullary strict node have both
\(\operatorname{Def}\) and \(\operatorname{Last}\) empty.
A defeasible top node has its rule in \(\operatorname{Def}\) and has the
nonempty last-rule set \(\{r\}\). At a strict top node,
\(\operatorname{Last}\) is empty exactly when every child's last-rule set
is empty. By induction this is equivalent to every child's defeasible-rule
set being empty, which is exactly emptiness of the parent's
\(\operatorname{Def}\). This proves the equivalence.
\end{proof}
\end{annexitem}

\begin{annexitem}{annex:p18.aspicsets}{thm:p18-aspic-set-properties}{Set comparisons and strict parts}
\begin{proof}
Suppose \(\Gamma\unlhd\Delta\) and \(\Delta\unlhd\Theta\).
The first relation forces \(\Gamma\ne\varnothing\), and the second forces
\(\Delta\ne\varnothing\). If \(\Theta=\varnothing\), the desired relation
\(\Gamma\unlhd\Theta\) follows from its empty-right-operand clause.
Assume henceforth that all three sets are nonempty.

For the elitist comparison choose \(x\in\Gamma\) with
\(x\leq y\) for every \(y\in\Delta\), and choose \(y_0\in\Delta\)
with \(y_0\leq z\) for every \(z\in\Theta\). Transitivity of the
base preorder gives \(x\leq z\) for every \(z\in\Theta\).
For the democratic comparison, for each \(x\in\Gamma\) choose
\(y\in\Delta\) with \(x\leq y\), and then \(z\in\Theta\)
with \(y\leq z\). Thus \(x\leq z\), proving the required
universal-existential condition. Both relations are transitive.

For any transitive relation \(R\), its asymmetric part
\(xPy\iff xRy\land\neg(yRx)\) is irreflexive. If \(xPy\) and \(yPz\),
then \(xRz\). Were \(zRx\), transitivity with \(yRz\) would give \(yRx\),
contrary to \(xPy\). Hence \(xPz\). Apply this argument to either set
comparison. No reflexivity of \(R\) was needed.
The empty-left-operand clause proves failure of reflexivity at the empty
set. For the elitist comparison it may also fail at a nonempty set of
pairwise incomparable elements.

In a nonempty finite set of a total preorder, a least element exists:
start with one element and successively retain the smaller of it and each
new element. The retained element is below every element considered so far,
by transitivity. Reversing this procedure produces a greatest element.
If \(g,d\) are least elements of \(\Gamma,\Delta\), respectively, then
\(\Gamma\unlhd_{\rm E}\Delta\) implies \(g\leq x\leq d\) for an
elitist witness \(x\). Conversely \(g\leq d\) implies \(g\leq y\)
for every \(y\in\Delta\), so \(g\) witnesses the relation.
If \(g,d\) instead denote greatest elements, the democratic relation
implies \(g\leq y\leq d\) for a witness \(y\in\Delta\), and
\(g\leq d\) implies that \(d\) witnesses the condition for every
\(x\in\Gamma\). This proves both extremal descriptions.

For the counterexample let the base preorder on \(\{a,b\}\) be equality.
Set \(\Gamma=\{a,b\}\), \(\Delta=\{a\}\).
The witness \(a\) proves \(\Gamma\unlhd_{\rm E}\Delta\).
The reverse relation would require \(a\leq b\), which is false.
Thus \(\Gamma\lhd_{\rm E}\Delta\).
The strict part of the base equality preorder is empty, however, so no
\(x\in\Gamma\) is strictly below every element of \(\Delta\).
The direct strict substitution therefore gives a different answer.

Even with total priorities, strict comparison of several components is not
generally obtained by requiring that each component be strictly smaller.
For example, take arguments with one ordinary premise and one defeasible
rule each, with their premise priorities equal and the first rule strictly
less preferred than the second. In the third weakest-link case the first
argument weakly compares below the second and the reverse comparison fails;
its strict asymmetric part holds. Requiring strict inequality of both
premise and rule comparisons would reject it. This explains why the body
defines strict argument preference only after the whole weak comparison.
\end{proof}
\end{annexitem}

\begin{annexitem}{annex:p18.aspiclocality}{thm:p18-aspic-locality}{Local targets and propagation}
\begin{proof}
Suppose an attack of \(A\) on \(B\) uses the target \(B'\).
For undermining, \(B'\) is the premise subtree \([\phi]\); for either
other mode, \(B'\) is the identified rule subtree. If
\(B\in\operatorname{Sub}(T)\), tree transitivity gives
\(B'\in\operatorname{Sub}(T)\).
The premise type, rule type, name, and conflict facts have not changed,
so the same attack targets \(T\).
An unconditional success remains unconditional. For a
preference-sensitive success the comparison remains
\(\neg(A\prec B')\), with exactly the same two arguments.
It therefore remains a defeat.

Now let \(T=[A_1,\ldots,A_k;r]\) have strict top rule.
Its root is not a premise node and is neither a defeasible rebuttal target
nor a defeasible name target. Every attack on \(T\) therefore targets a
proper subargument contained in some \(A_i\). Its target and preference
comparison are identical in that child, so any successful attack defeats
that \(A_i\). A nullary strict node has no proper subarguments, and
therefore has no attacks or defeats at all.
\end{proof}
\end{annexitem}

\begin{annexitem}{annex:p18.aspicqueries}{thm:p18-aspic-query-distinction}{Formula witnesses}
\begin{proof}
If one argument \(A\) with conclusion \(\phi\) belongs to every extension,
that same \(A\) supplies the existential witness in each extension.
The nonempty-family guard is identical for both queries.

For the converse, let \(L=\{a,b,p,\bar p\}\), with
\(C(a)=\{b\}, C(b)=\{a\}, C(p)=\{\bar p\}, C(\bar p)=\{p\}\).
Let \(K_p=\{a,b\}\), \(K_n=R_s=\varnothing\), and let the two
unnamed defeasible rules be \(r_a:a\Rightarrow p\) and
\(r_b:b\Rightarrow p\). Use the empty strict preference relation.
Exactly four arguments exist:
\[
 A=[a],\quad B=[b],\quad P=[A;r_a],\quad Q=[B;r_b].
\]
Their defeat edges are exactly
\[
 A D B,\quad A D Q,\quad B D A,\quad B D P.
\]
The first pair comes from undermining \(b\), and the second pair from
undermining \(a\). There is no argument concluding \(\bar p\), and no
rule is named, so there are no additional rebuttals or undercuts.

The sets \(\{A,P\}\) and \(\{B,Q\}\) are conflict-free and attack
every argument outside themselves; thus both are stable.
An admissible set containing \(P\) must defeat its attacker \(B\).
Only \(A\) defeats \(B\), so it must contain \(A\).
Similarly, containing \(Q\) requires containing \(B\).
Conflict-freeness prevents simultaneous membership of \(A\) and \(B\).
It follows that every nonempty admissible set is contained in one of
\(\{A,P\},\{B,Q\}\); the empty set is not maximal. These are therefore
exactly the preferred extensions and exactly the stable extensions.
Each contains an argument for \(p\), but their intersection is empty.
Thus \(p\) is extensionwise skeptical and is not common-argument skeptical
for either of these two extension families.
\end{proof}
\end{annexitem}

\begin{annexitem}{annex:p18.aspicclosure}{thm:p18-aspic-structural-rationality}{Complete extensions contain subarguments and strict extensions}
\begin{proof}
First justify the closure definition. Its iterates are increasing.
If every antecedent of a finite-arity strict rule belongs to
\(\bigcup_jS_j\), there is a single \(m\) with all its antecedents in
\(S_m\): take the maximum of their finitely many stages. Its consequent
belongs to \(S_{m+1}\). A nullary rule has its consequent in \(S_1\).
Thus the union is closed under all strict rules.
If \(T\) contains \(S\) and is closed under these rules, induction gives
\(S_j\subseteq T\) for every \(j\). Consequently the union is the least
such closed superset.

Let \(E\) be complete, \(A\in E\), and \(B\in\operatorname{Sub}(A)\).
Every defeater \(C\) of \(B\) defeats \(A\) by
Lemma~\ref{thm:p18-aspic-locality}.
Since \(E\) is admissible and contains \(A\), it contains a defeater of
every such \(C\). Hence \(E\) defends \(B\).
Completeness includes every defended argument, so \(B\in E\).
This proves subargument closure without an additional consistency premise.

Next suppose a strict rule \(r\) has antecedents
\(\phi_1,\ldots,\phi_k\) in \(\operatorname{Conc}(E)\).
Choose \(A_i\in E\) with conclusion \(\phi_i\), and construct
\(T=[A_1,\ldots,A_k;r]\). Such a tree belongs to the unrestricted
argument set by definition.
Every defeater of \(T\) defeats one of the \(A_i\), by the same locality
lemma. Admissibility of \(E\) therefore provides a counter-defeater in
\(E\). Thus \(E\) defends \(T\), and completeness gives \(T\in E\).
For \(k=0\), \(T\) has no defeaters, so the same conclusion holds by
vacuous defense. It follows that \(\operatorname{Conc}(E)\) is
strictly closed; the least-closure property proved above gives the
displayed equality.

The argument used neither finiteness of \(\mathcal A\) nor a property of
the preference relation beyond its fixed local use in defeat.
The already established abstract extension relations identify grounded,
preferred, and stable extensions as complete, so the same conclusions
apply to each of them.
\end{proof}
\end{annexitem}

\begin{annexitem}{annex:p18.aspicflat}{thm:p18-aspic-flat-consistency}{Consistency with no strict rules or axioms}
\begin{proof}
With \(K_n=R_s=\varnothing\), every argument is either an ordinary
premise leaf or has a defeasible top rule. Thus its conclusion is always
an undermining or rebutting target at its root.
Suppose a conflict-free set \(E\) had conclusions
\(\alpha=\operatorname{Conc}(A)\), \(\beta=\operatorname{Conc}(B)\)
with \(\alpha\in C(\beta)\), for \(A,B\in E\).
If \(\beta\notin C(\alpha)\), \(A\)'s root-directed attack on \(B\)
is contrary and succeeds independently of preference.
This gives \(A D B\), contrary to conflict-freeness.

If also \(\beta\in C(\alpha)\), both root-directed attacks exist.
Preventing \(A D B\) would require \(A\prec B\), while preventing
\(B D A\) would require \(B\prec A\). Asymmetry excludes their conjunction,
so at least one defeat edge exists within \(E\), again a contradiction.
This reasoning also covers \(A=B\): asymmetry implies \(\neg(A\prec A)\),
and a self-conflicting conclusion then supplies a self-defeat.

The conclusion set is therefore directly consistent. Since \(R_s\) is
empty its strict closure equals itself, so it is indirectly consistent
as well. For a complete \(E\), the structural closure theorem supplies
subargument closure and strict closure. These are the four asserted
properties in this stated, restricted domain.
\end{proof}
\end{annexitem}

\begin{annexitem}{annex:p18.aspicboundaries}{thm:p18-aspic-rationality-boundaries}{Complete countermodels}
\begin{proof}
All examples use paired literals, with each positive literal conflicting
only with its negative partner, and conversely. Unmentioned premise sets,
rule sets, and names are empty. Strict preferences are empty.

\emph{Conflicting axioms.}
Take \(L=\{p,\neg p\}\), \(K_n=\{p,\neg p\}\), and \(K_p=\varnothing\).
There are exactly two arguments, \([p]\) and \([\neg p]\).
Neither has an ordinary premise or a defeasible top rule or named
defeasible rule, so neither is attackable. The induced framework has no
edges. Both arguments are defended by every set; consequently its only
complete extension is the entire argument set. That set is also grounded,
preferred, and stable. Its conclusions are inconsistent. No choice of
preferences can remove this defect.

\emph{Strict rules without the needed transpositions.}
Use paired literals for \(a,b,p\), ordinary premises \(a,b\), and strict
rules \(r:a\to p\), \(s:b\to\neg p\). There are exactly
\[
 A=[a],\quad B=[b],\quad P=[A;r],\quad N=[B;s].
\]
The premise set \(\{a,b\}\) is directly consistent.
The only conflicting conclusions are \(p,\neg p\), and both are at
strict roots. Neither \(\neg a\) nor \(\neg b\) is concluded.
Hence there are no attacks: all four arguments form the unique extension,
which concludes both \(p\) and \(\neg p\).
The transposed rules \(\neg p\to\neg a\) and \(p\to\neg b\)
are absent, explaining the explicit closure failure in the rule data.

One can calculate the effect of adding exactly these two rules without
appealing to a general rationality theorem. They add just
\[
 U=[P;p\to\neg b],\qquad V=[N;\neg p\to\neg a].
\]
The complete list of defeats is
\[
 U D B,\ U D N,\ U D V,\qquad
 V D A,\ V D P,\ V D U .
\]
The two sets \(\{A,P,U\}\) and \(\{B,N,V\}\) are stable.
Every admissible set that contains a member of the first group must
defeat \(V\), and only \(U\) does so; every admissible set that contains
a member of the second must contain \(V\).
The two groups cannot coexist because \(U,V\) defeat each other.
Thus these are exactly the preferred and stable extensions.
The empty set is grounded because every argument has an attacker.
Each displayed extension now has a consistent conclusion set.

\emph{Allowing rebuttal at strict roots.}
Use paired literals for \(a,b,p,q\), \(K_p=\{a,b\}\), and rules
\[
 d:a\Rightarrow p,\quad e:b\Rightarrow q,\qquad
 r:p\to\neg q,\quad s:q\to\neg p .
\]
The two strict rules are each other's transposition under the paired-literal
involution. Exactly six arguments exist:
\[
 A_0=[a],\quad B_0=[b],\quad
 A=[A_0;d],\quad B=[B_0;e],\quad
 S=[A;r],\quad T=[B;s].
\]
Under the chapter's actual attack definition, the edges are
\[
 S D B,\ S D T,\qquad T D A,\ T D S .
\]
Now change only rebuttal: also permit a conflicting attack at a strict
rule root. This adds \(A D T\) and \(B D S\); it adds no further edges.
The set \(E=\{A_0,B_0,A,B\}\) has no internal edge and defeats both
outsiders \(S,T\), through \(B D S\) and \(A D T\).
It is therefore stable, and hence admissible and preferred, in this
modified framework.
Its conclusions are \(\{a,b,p,q\}\), whereas strict inference adds
\(\neg q,\neg p\). Thus it is not strictly closed and its strict closure
is inconsistent. This countermodel uses an explicitly modified attack
definition; it does not contradict the structural theorem for the
target-local ASPIC+ construction.
\end{proof}
\end{annexitem}

\begin{annexitem}{annex:p18.aspicgraph}{thm:p18-aspic-worked-graph}{Exhaustive running example and its extensions}
\begin{proof}
The only available premise conclusions are \(p,s,u,x\), each with its
single listed premise argument. The only rule producing \(q\) is \(d_1\)
from \(p\); the only rule producing \(r\) is \(s_1\) from \(p,q\).
Similarly, \(t\) comes only from \(s\) by \(d_2\), then \(\neg d_1\)
comes only from \(t\) by \(d_3\), and \(\neg p\) comes only from \(s\)
by \(d_6\). The conclusion \(v\) comes only from \(u\) by \(d_4\);
it produces \(\neg s\) by \(s_2\), and, together with the unique \(x\)
argument, produces \(\neg t\) by \(d_5\).
No other rule consequent feeds an additional rule antecedent.
This proves that the twelve listed trees, and no others, are generated.
Their component sets follow directly from the recursive definitions.

For undermining, the only ordinary-premise targets have conclusions
\(s,u,x\). Of their contradictories, only \(\neg s\) is concluded:
by \(C_3\). The subargument \([s]=B_1\) occurs exactly in
\(B_1,B_2,B_3,B_4\), giving all four listed undermining edges.
The contradictory \(\neg p\) of the axiom premise \(p\) provides no
undermining attack.

For rebuttal, the possible target-root conclusions are
\(q,t,\neg d_1,\neg p,v,\neg t\).
Among their contradictories, only \(p,t,\neg t\) occur as argument
conclusions. The conclusion \(p\), supplied by \(A_1\), contradicts
the defeasible conclusion \(\neg p\) of \(B_4\), so there is
the rebuttal \(A_1 D B_4\) when it succeeds.
The arguments \(D_4\) and \(B_2\) yield the two directions between
\(\neg t\) and \(t\), with the attack on \(B_2\) propagating to
\(B_3\). No argument concludes \(d_1,\neg q,\neg v\).
For undercutting, the only negative name conclusion is
\(\neg d_1\), from \(B_3\), and rule \(d_1\) occurs exactly in
\(A_2,A_3\). These give the two undercuts.

Thus the complete structural attack list is the body table, including its
\(A_1\)-to-\(B_4\) rebuttal. With no strict preferences every attack
succeeds. The initially unattacked arguments are exactly
\[
 U_0=\{A_1,C_1,C_2,C_3,D_3\}.
\]
They belong to every complete extension. Since \(C_3\) defeats every
\(B_i\), no \(B_i\) can belong to an admissible set: defending it would
require defeating the unattacked \(C_3\). Also \(A_1\) independently
defeats \(B_4\).
The set \(U_0\) defeats the only attackers of \(A_2,A_3,D_4\):
their attackers are \(B_3,B_3,B_2\), respectively.
Every complete extension must therefore include these three as well.
The resulting \(E_0\) is conflict-free and defeats every outsider, namely
the four \(B_i\). It is stable and hence complete.
The forced inclusions and exclusions prove uniqueness among complete
extensions. Grounded and preferred extensions exist and are complete;
stable extensions are complete; consequently the same \(E_0\) is unique
in all four semantics. Finally, \(A_3\in E_0\) concludes \(r\), so both
formula-skepticism conventions hold for each singleton extension family.
\end{proof}
\end{annexitem}

\begin{annexitem}{annex:p18.aspicpreferences}{thm:p18-aspic-worked-preferences}{Computing the link comparisons}
\begin{proof}
All conflict comparisons are symmetric in this example, and undercuts
remain unconditional. The only preference-sensitive pairs are
\((C_3,B_1)\), \((D_4,B_2)\), \((B_2,D_4)\), and \((A_1,B_4)\).

For last link,
\(\operatorname{Last}(C_3)=\{d_4\}\) and
\(\operatorname{Last}(B_1)=\varnothing\).
The empty-set clauses give \(C_3\prec_{\rm L}B_1\), so none of the
four undermining attacks succeeds.
Also \(\operatorname{Last}(B_2)=\{d_2\}\),
\(\operatorname{Last}(D_4)=\{d_5\}\), and \(d_2<_d d_5\).
Hence \(B_2\prec_{\rm L}D_4\), while the reverse strict comparison
is false. The successful rebuttals are \(D_4 D B_2,D_4 D B_3\).
Finally \(A_1\) has empty last-rule and ordinary-premise sets.
Since \(B_4\) has last rule \(d_6\), it is strictly below \(A_1\).
Thus \(A_1\) successfully rebuts \(B_4\), regardless of the rank of \(d_6\).
The complete defeat graph is therefore
\[
 B_3D A_2,\ B_3D A_3,\quad
 D_4D B_2,\ D_4D B_3,\quad A_1D B_4 .
\]
Its unattacked vertices are
\(\{A_1,B_1,C_1,C_2,C_3,D_3,D_4\}\).
Every complete extension includes them, excludes \(B_2,B_3,B_4\),
and includes \(A_2,A_3\), since \(D_4\) defeats their sole attacker
\(B_3\). The resulting set is conflict-free and attacks all three
outsiders. It is stable, and the forced choices prove uniqueness in all
four semantics.

For weakest link, \(C_3\) is plausible and defeasible whereas \(B_1\)
is plausible and strict. The third case applies:
\(\{u\}\lhd_{\rm E}\{s\}\) and
\(\{d_4\}\lhd_{\rm E}\varnothing\).
Thus \(C_3\prec_{\rm W}B_1\), and undermining again fails.
Both \(D_4,B_2\) are plausible and defeasible.
The ordinary-premise sets satisfy
\[
 \{u,x\}\unlhd_{\rm E}\{s\},\qquad
 \neg(\{s\}\unlhd_{\rm E}\{u,x\}),
\]
because both \(u,x\) are strictly less preferred than \(s\).
Their defeasible-rule sets satisfy
\[
 \{d_4,d_5\}\unlhd_{\rm E}\{d_2\},\qquad
 \neg(\{d_2\}\unlhd_{\rm E}\{d_4,d_5\}),
\]
because \(d_4<_d d_2<_d d_5\).
Consequently \(D_4\prec_{\rm W}B_2\), so precisely the rebuttal
\(B_2 D D_4\) succeeds in this pair.
Comparing \(B_4\) with the strict firm argument \(A_1\), the third
case gives both \(\{s\}\unlhd\varnothing\) and
\(\{d_6\}\unlhd\varnothing\); the reverse comparisons fail.
Thus \(B_4\prec_{\rm W}A_1\) and \(A_1 D B_4\).
The complete graph is
\[
 B_3D A_2,\ B_3D A_3,\qquad B_2D D_4,\qquad A_1D B_4 .
\]
Every other argument is unattacked, so every complete extension includes
exactly those other arguments and excludes these four attacked vertices.
That set is conflict-free and stable; uniqueness follows as before.
It contains no \(A_3\), and no other argument concludes \(r\).

In both cases \(B_1,C_3\) are included and conclude \(s,\neg s\).
Both graphs exclude \(B_4\), so neither extension contains an argument
concluding \(\neg p\). The pair \(s,\neg s\) proves the stated
inconsistency boundary.
\end{proof}
\end{annexitem}

\begin{annexitem}{annex:p18.aspiccontrary}{thm:p18-aspic-contrary-example}{Contrary defeat ignores relative strength}
\begin{proof}
There are exactly two arguments \(P=[p]\) and \(S=[s]\).
The pairs \(p,\bar p\) and \(s,\bar s\) are symmetric conflicts, so
each formula has a contradictory. The additional conflict satisfies
\(p\in C(s)\) and \(s\notin C(p)\), hence it is contrary.
The argument \(P\) therefore contrary-undermines \(S\).
There is no attack in the reverse direction and neither argument
self-attacks.
Both arguments are strict and have singleton ordinary-premise sets.
For either link principle and either set lift, \(p<_p s\) yields
\(P\prec S\). Contrary undermining succeeds nevertheless.
The sole defeat edge is \(P D S\). Every complete extension includes
the unattacked \(P\), and no admissible set contains \(S\), whose
attacker \(P\) cannot be defeated. Thus \(\{P\}\) is the unique complete
extension and is stable; the four semantics coincide.
\end{proof}
\end{annexitem}

\begin{annexitem}{annex:p18.aspicfinite}{thm:p18-aspic-finite-generation}{A terminating generation method and its boundary}
\begin{proof}
There are finitely many formulas occurring in the finite knowledge base
and the finitely many finite-arity rules. The rule-name formulas and
their conflict data need not generate arguments unless a premise or
rule consequent supplies them.
For each consequent formula of rank zero, the only arguments are its
premise leaf, if present, and its finitely many nullary-rule nodes:
a non-nullary rule cannot have an antecedent of smaller nonnegative rank.
Thus there are finitely many arguments with rank-zero conclusions.

Suppose all argument sets with conclusion rank below \(j\) have already
been generated and are finite. Every rule with consequent rank \(j\)
has either no antecedents or only antecedents of smaller rank.
For antecedents \((\phi_1,\ldots,\phi_k)\), its possible children form
the finite Cartesian product of the already generated argument sets
concluding \(\phi_i\). Generate one new labelled ordered tree for each
tuple. Repeated antecedent formulas still occupy their specified
positions, so this generates all ordered choices.
There are finitely many rules and premise leaves at rank \(j\);
therefore this stage is finite. An empty child-argument set yields no
tuple for that non-nullary rule. Nullary products have exactly one
empty tuple. Induction through \(j=H\) proves termination.

To see exhaustiveness, consider any argument tree. If it is a premise
or nullary node, it was generated. Otherwise its children's conclusions
have strictly lower rank than its conclusion, so by induction on rank
all its children had been generated before the rule's stage; their
ordered tuple produces this tree. Soundness is immediate because every
generated tuple obeys the rule's antecedent labels. Duplicate records
of the identical labelled tree may simply be identified.

Once the finite list is present, recursion computes each argument's
subarguments, premises, last rules, and all defeasible rules. For each
ordered argument pair, inspect each possible target subargument and
check the appropriate conflict membership and name condition. Each set
comparison quantifies only over finite component sets with supplied
priority data. Thus every defeat edge is determined by finitely many
given comparisons. Apply the established finite abstract methods to this
explicit graph; formula queries then range over finite extension and
argument lists. This is a decidability statement for explicitly supplied
data, not a claim about an oracle for arbitrary infinite language relations.

For the boundary, let \(A_0=[p]\) and recursively
\(A_{m+1}=[A_m;r]\), where \(r:p\to p\).
Each \(A_m\) is an argument of height \(m\). Distinct heights imply
distinct trees, so \(\{A_m:m\geq0\}\) is infinite despite the single
premise and single rule. Identifying all these trees by their common
conclusion would change the structured argument set and is not the
construction proved above.
\end{proof}
\end{annexitem}
